\begin{abstract}
LLM agents act over many turns (tool calls, searches, memory writes) but are usually trained with a single outcome reward, so reinforcement learning must decide which turns deserve the credit. GRPO and its 2025--2026 variants answer ``all of them, equally'': they refine how a uniform per-response advantage is optimized, not who receives it. We survey credit assignment for LLM agents, focusing on 2025--2026 work. We organize methods into six families (process rewards, Monte-Carlo and tree branching, anchor-state grouping, turn-level rewards and critics, hindsight and counterfactual credit, revived value models) and profile each by what it requires of the environment and where its extra compute goes. We read Memory-R2 closely: its LoGo-GRPO adds rerollouts from a shared memory state to fix unfair group baselines, but its gain over GRPO lacks a compute-matched control. A pre-registered study of turn-level credit in math tool use gives three lessons: non-conservative credit acts as a reward and is gamed (over-calling); conservation removes the pathology but yields no detectable gain at four turns or fewer; and DAPO-style filtering's endpoint gain, bought with 2.53$\times$ the rollouts, vanishes at matched compute. The field has many methods but almost no way to measure whether credit is correct. Our primary recommendation is a controllable-horizon benchmark with computable ground-truth credit, linked to compute-matched training.
\end{abstract}
